\section{Inledning}
\label{sec:inledning}

Jag har gjort min andra och sista LIA-period på Insighta, Inc, ett företag inom marketing
analytics med kontor i Knoxville, USA. Företaget hjälper sina kunder att samla in, köra och
analysera data så att de kan fatta bättre beslut. Min roll under perioden var DevOps Engineer
och min handledare var Amir Jaber. LIA-perioden pågick från januari till
maj 2026.

\subsection*{Varför jag valde samma företag igen}

Jag valde att göra LIA 2 på samma företag som LIA 1 eftersom jag ville bygga vidare på det
jag redan hade påbörjat. Under LIA 1 fick jag en bra grund i företagets arbetssätt och teknik,
men jag hann inte arbeta tillräckligt mycket med drift i produktion. Jag ville därför använda
LIA 2 till att förstå hur systemen fungerar när de används på riktigt och när fel faktiskt får
konsekvenser.

Jag funderade också på att välja ett större företag. Det hade kunnat ge en tydligare
organisationsstruktur och mer specialiserade roller. Men jag kom fram till att ett mindre team
passade mina mål bättre. På Insighta fick jag arbeta nära hela kedjan: från Terraform och
GitHub Actions till Kubernetes, säkerhet och driftproblem. Det gav mig mer ansvar och mer
lärande.

\subsection*{Mina mål med LIA 2}

LIA 2 hade tydligt fokus på drift, stabilitet och automatisering. Mina mål var att:

\begin{itemize}
  \item flytta företagets datapipelinesystem från Docker på virtuella maskiner till ett
        Kubernetes-kluster i Google Cloud,
  \item använda Helm och Terraform för att göra deployment mer strukturerad och repeterbar,
  \item införa Workload Identity så att systemet kan använda säkrare autentisering utan
        statiska nycklar,
  \item automatisera deployment med GitHub Actions,
  \item minska driftstopp när ny kod deployas,
  \item hålla CI/CD-pipelinen uppdaterad när verktyg och versioner förändras.
\end{itemize}

Dessa mål hänger tydligt ihop med det vi har läst i utbildningen, särskilt inom molntjänster,
nätverk, operativsystem och IT-säkerhet. Skillnaden under LIA var att jag behövde använda
kunskaperna i en verklig miljö där lösningarna måste fungera för teamet och för företagets
kunder.